\documentclass[11pt]{ctexart}
\usepackage[a4paper,margin=1in]{geometry}
\usepackage{amsmath,amsfonts,amssymb}
\usepackage{graphicx}
\usepackage{booktabs}
\usepackage{caption}
\usepackage{enumitem}
\usepackage[hidelinks]{hyperref}
\usepackage{microtype}

\title{Lazy 环 + 单向 shortcut：$\beta$ 扫描、尾部机制与 jump-over 扩展报告}
\author{}
\date{}

\setlist{nosep}
\setlength{\textfloatsep}{8pt plus 2pt minus 2pt}
\setlength{\floatsep}{6pt plus 2pt minus 2pt}
\setlength{\intextsep}{6pt plus 2pt minus 2pt}
\setlength{\abovecaptionskip}{4pt}
\setlength{\belowcaptionskip}{0pt}

\begin{document}
\maketitle

\section*{结论与关键图示}
固定 $N=100,q=2/3,\rho=1$ 且仅改变 $\beta$ 时，
两峰整体向更早时间移动，尾部衰减率 $\gamma$ 随 $\beta$ 增大而上升；
同时 $K=4$ 的双峰区间更宽、valley 更深。图~\ref{fig:beta-sweep-key-cn} 汇总了峰高、峰位与尾部指标的变化。

\begin{figure}[!htbp]
  \centering
  \begin{minipage}{0.48\linewidth}
    \centering
    \textbf{(a) 峰高}\par
    \includegraphics[width=\linewidth]{outputs/beta_sweep_N100/peak_heights_vs_beta_N100.png}
  \end{minipage}\hfill
  \begin{minipage}{0.48\linewidth}
    \centering
    \textbf{(b) 峰位}\par
    \includegraphics[width=\linewidth]{outputs/beta_sweep_N100/peak_times_vs_beta_N100.png}
  \end{minipage}
  \vspace{4pt}
  \begin{minipage}{0.48\linewidth}
    \centering
    \textbf{(c) 尾部 $\gamma$}\par
    \includegraphics[width=\linewidth]{outputs/beta_sweep_N100/tail_gamma_vs_beta_N100.png}
  \end{minipage}\hfill
  \begin{minipage}{0.48\linewidth}
    \centering
    \textbf{(d) 全 $\beta$ AW 叠加}\par
    \includegraphics[width=\linewidth]{outputs/beta_sweep_N100/aw_all_betas_N100.png}
  \end{minipage}
  \caption{仅改变 $\beta$ 的精确 AW 结果：峰高、峰位、tail 衰减率 $\gamma$ 与全 $\beta$ 叠加的 $f(t)$；$\beta=0$ 为无 shortcut 的基准点，单独标记。}
  \label{fig:beta-sweep-key-cn}
\end{figure}

以 $\beta=0.01$ 为例，代表性 $f(t)$ 更直观地展示 valley 深度差异：$K=2$ 的
$\texttt{hv\_over\_max}\approx0.34$，而 $K=4$ 约 $0.15$，说明 $K=4$ 两峰分离更明显；
同一 $\beta$ 下 $K=4$ 的 $(t_1,t_2)$ 更早且 $\gamma$ 更大，时间尺度更快、尾更轻。
图~\ref{fig:beta-example-ft-cn} 给出对应的峰谷窗口。

\begin{figure}[!htbp]
  \centering
  \begin{minipage}{0.48\linewidth}
    \centering
    \textbf{(a) $K=2$}\par
    \includegraphics[width=\linewidth]{outputs/mc_beta_sweep_N100/cases/N100_K2_beta0.01/N100_K2_beta0.f_t.pdf}
  \end{minipage}\hfill
  \begin{minipage}{0.48\linewidth}
    \centering
    \textbf{(b) $K=4$}\par
    \includegraphics[width=\linewidth]{outputs/mc_beta_sweep_N100/cases/N100_K4_beta0.01/N100_K4_beta0.f_t.pdf}
  \end{minipage}
  \caption{代表性 $f(t)$（AW 精确，已聚焦峰谷时间窗）。}
  \label{fig:beta-example-ft-cn}
\end{figure}

尾部机制由 transient 谱半径控制，$\gamma=-\log(\lambda_{\max})$；
$\log S(t)$ 在尾部近线性衰减，与谱半径指标一致（图~\ref{fig:tail-logS-cn}）。

\begin{figure}[!htbp]
  \centering
  \begin{minipage}{0.48\linewidth}
    \centering
    \textbf{(a) $K=2$}\par
    \includegraphics[width=\linewidth]{outputs/tail_diag/log_survival_N100_K2_beta0.01.png}
  \end{minipage}\hfill
  \begin{minipage}{0.48\linewidth}
    \centering
    \textbf{(b) $K=4$}\par
    \includegraphics[width=\linewidth]{outputs/tail_diag/log_survival_N100_K4_beta0.01.png}
  \end{minipage}
  \caption{生存函数 $S(t)$ 的对数衰减（$\beta=0.01$），与谱半径指标 $\gamma$ 对应。}
  \label{fig:tail-logS-cn}
\end{figure}

MC 轨迹分类进一步揭示 jump-over 只在 $K=4$ 出现：
$K=2$ 的 $J$ 类几乎为 0，peak1 基本全是 $\texttt{C1pJ0}$，
valley 以 $\texttt{C1pJ0}$ 为主并夹少量 $\texttt{C0J0}$，peak2 则更多落在 $\texttt{C0J0}/\texttt{C1pJ0}$ 上。
$K=4$ 的结构更清晰：peak1 主要由 $\texttt{C1pJ0}$ 与 $\texttt{C1pJ1p}$ 构成，$\texttt{C1pJ1p}$ 约占三分之一；
valley 中 $\texttt{C1pJ1p}$ 也在三分之一左右，显示 jump-over 对 valley 延迟路径的稳定贡献；
peak2 则被 $\texttt{C0}$ 类主导，$\texttt{C0J0}$ 为主并伴随约四分之一的 $\texttt{C0J1p}$，
这些 $J1p$ 轨迹对应“跨过吸收点但未吸收”的延迟路径。
从图中可以看出这类比例随 $\beta$ 仅小幅起伏（图~\ref{fig:mc-beta-classes-cn}）。
这里回顾四类分类的含义：$C$ 表示是否至少使用一次 shortcut（$C0$ 为未用、$C1p$ 为至少一次），
$J$ 表示是否出现 jump-over（$J0$ 为无、$J1p$ 为至少一次），
因此 $\texttt{C0J0}$、$\texttt{C1pJ0}$、$\texttt{C0J1p}$、$\texttt{C1pJ1p}$ 分别对应“无 shortcut/无 jump-over”、
“有 shortcut/无 jump-over”、“无 shortcut/有 jump-over”、“有 shortcut/有 jump-over”的轨迹类型。

\begin{figure}[!htbp]
  \centering
  \begin{minipage}{0.32\linewidth}
    \centering
    \textbf{(a) peak1}\par
    \includegraphics[width=\linewidth]{outputs/mc_beta_sweep_N100/classes_peak1_vs_beta.png}
  \end{minipage}\hfill
  \begin{minipage}{0.32\linewidth}
    \centering
    \textbf{(b) valley}\par
    \includegraphics[width=\linewidth]{outputs/mc_beta_sweep_N100/classes_valley_vs_beta.png}
  \end{minipage}\hfill
  \begin{minipage}{0.32\linewidth}
    \centering
    \textbf{(c) peak2}\par
    \includegraphics[width=\linewidth]{outputs/mc_beta_sweep_N100/classes_peak2_vs_beta.png}
  \end{minipage}
  \caption{四类轨迹在不同窗口内的条件概率随 $\beta$ 的变化（$N=100$）。}
  \label{fig:mc-beta-classes-cn}
\end{figure}

为选取代表性双峰参数，采用统一阈值 $t_2/t_1\ge10$、$\texttt{hv\_over\_max}\le0.6$、
$\texttt{h2\_over\_h1}\ge0.05$，取同时满足 $K=2$ 与 $K=4$ 的最小 $\beta$，得到 $\beta_\ast=0.01$。

固定 $\beta_\ast=0.01$ 扫 $N$ 时，$t_2$ 随 $N$ 增大显著右移、$h_2$ 下降、$\gamma$ 下降（尾部更厚），
而 $t_1$ 基本稳定；图~\ref{fig:n-sweep-peaks-cn} 与图~\ref{fig:n-sweep-gamma-cn} 给出定量结果。

\begin{figure}[!htbp]
  \centering
  \begin{minipage}{0.48\linewidth}
    \centering
    \textbf{(a) $t_1$}\par
    \includegraphics[width=\linewidth]{outputs/N_sweep_beta_selected/t1_vs_N_beta0.01.png}
  \end{minipage}\hfill
  \begin{minipage}{0.48\linewidth}
    \centering
    \textbf{(b) $t_2$}\par
    \includegraphics[width=\linewidth]{outputs/N_sweep_beta_selected/t2_vs_N_beta0.01.png}
  \end{minipage}
  \vspace{4pt}
  \begin{minipage}{0.48\linewidth}
    \centering
    \textbf{(c) $h_1$}\par
    \includegraphics[width=\linewidth]{outputs/N_sweep_beta_selected/h1_vs_N_beta0.01.png}
  \end{minipage}\hfill
  \begin{minipage}{0.48\linewidth}
    \centering
    \textbf{(d) $h_2$}\par
    \includegraphics[width=\linewidth]{outputs/N_sweep_beta_selected/h2_vs_N_beta0.01.png}
  \end{minipage}
  \caption{固定 $\beta=0.01$ 的峰位/峰高随 $N$ 的变化。}
  \label{fig:n-sweep-peaks-cn}
\end{figure}

\begin{figure}[!htbp]
  \centering
  \includegraphics[width=0.7\linewidth]{outputs/N_sweep_beta_selected/tail_gamma_vs_N_beta0.01.png}
  \caption{固定 $\beta=0.01$ 时 tail 衰减率 $\gamma$ 随 $N$ 的变化。}
  \label{fig:n-sweep-gamma-cn}
\end{figure}

对应的 MC 扫 $N$ 进一步把结构拉直：$K=2$ 仍完全是 $J0$ 类，peak1 始终为 $\texttt{C1pJ0}$，
valley 以 $\texttt{C1pJ0}$ 为主，peak2 在 $\texttt{C0J0}$ 与 $\texttt{C1pJ0}$ 之间摆动；
$K=4$ 在 $N$ 较大、valley/第二峰清晰后，peak1 与 valley 的 $J1p$ 份额稳定在约三分之一，
peak2 主要为 $\texttt{C0J0}$ 并保留约四分之一的 $\texttt{C0J1p}$，显示 jump-over 的贡献随 $N$ 变化不大但持续存在（图~\ref{fig:mc-n-classes-cn}）。

\begin{figure}[!htbp]
  \centering
  \begin{minipage}{0.32\linewidth}
    \centering
    \textbf{(a) peak1}\par
    \includegraphics[width=\linewidth]{outputs/mc_N_sweep_beta_selected/classes_peak1_vs_N.png}
  \end{minipage}\hfill
  \begin{minipage}{0.32\linewidth}
    \centering
    \textbf{(b) valley}\par
    \includegraphics[width=\linewidth]{outputs/mc_N_sweep_beta_selected/classes_valley_vs_N.png}
  \end{minipage}\hfill
  \begin{minipage}{0.32\linewidth}
    \centering
    \textbf{(c) peak2}\par
    \includegraphics[width=\linewidth]{outputs/mc_N_sweep_beta_selected/classes_peak2_vs_N.png}
  \end{minipage}
  \caption{四类轨迹比例随 $N$ 变化（固定 $\beta=0.01$）。}
  \label{fig:mc-n-classes-cn}
\end{figure}

下文依次给出理论推导与 valley 定义，再展开 $\beta$ 扫描、MC 分类、$N$ 扫描、MC $N$ 与复现步骤。

\section{理论推导：生成函数、AW 反演与尾部衰减}

本节给出本文数值计算与尾部指标所依赖的核心理论推导。我们考虑在 $N$ 点环上、每步以概率 $q$ 做 $K=2k$ 邻居环走、以概率 $1-q$ 懒惰自环的随机游走，并在源点 $\mathrm{src}$ 增加一条到 $\mathrm{dst}$ 的有向 shortcut。

\subsection{无缺陷（无 shortcut）的谱分解与 resolvent}
记无缺陷转移矩阵为 $P_0$。对于非源点，单步转移为
\[
P_0(i\to i)=1-q,\qquad
P_0(i\to i\pm r)=\frac{q}{K},\quad r=1,\dots,k.
\]
由于 $P_0$ 平移不变（循环矩阵），其特征向量为 Fourier 模式
\[
\phi_\ell(x)=\exp\Big(\frac{2\pi i\ell x}{N}\Big),\qquad \ell=0,\dots,N-1,
\]
对应特征值为
\begin{equation}
\lambda_\ell
=(1-q)+\frac{q}{K}\sum_{r=1}^{k}\Big(e^{2\pi i\ell r/N}+e^{-2\pi i\ell r/N}\Big)
=(1-q)+\frac{2q}{K}\sum_{r=1}^{k}\cos\Big(\frac{2\pi \ell r}{N}\Big).
\end{equation}

定义 resolvent（生成函数矩阵）
\[
S^{(0)}(z)=(I-zP_0)^{-1},\qquad |z|<1.
\]
利用谱分解可得矩阵元
\begin{equation}
S^{(0)}(z)_{x,y}
=\frac{1}{N}\sum_{\ell=0}^{N-1}
\frac{\exp\Big(\frac{2\pi i\ell(x-y)}{N}\Big)}{1-z\lambda_\ell}.
\end{equation}
进一步可定义位移形式 $Q(d,z):=S^{(0)}(z)_{y+d,y}$，用于快速构造带缺陷情形。

\subsection{Selfloop 规则下的单向 shortcut：rank-one 更新}
在 \texttt{lazy\_selfloop} 模式下，shortcut 概率
\[
p=\beta(1-q),
\]
并从源点 $\mathrm{src}$ 的自环概率中扣除同等质量：$P(\mathrm{src}\to\mathrm{src})=(1-q)-p$，同时增加 $P(\mathrm{src}\to\mathrm{dst})=p$。因此转移矩阵可写为对 $P_0$ 的 rank-one 更新：
\begin{equation}
P = P_0 + p\, e_{\mathrm{src}}(e_{\mathrm{dst}}-e_{\mathrm{src}})^\top,
\end{equation}
其中 $e_i$ 为标准基向量。于是
\[
I-zP = (I-zP_0) - z p\, e_{\mathrm{src}}(e_{\mathrm{dst}}-e_{\mathrm{src}})^\top.
\]
应用 Sherman--Morrison 公式：
\[
(A+uv^\top)^{-1}=A^{-1}-\frac{A^{-1}uv^\top A^{-1}}{1+v^\top A^{-1}u},
\]
可在已知无缺陷 resolvent $S^{(0)}(z)$ 的基础上，以 $O(1)$ 次矩阵元操作得到带缺陷 resolvent $S(z)$ 的所需矩阵元（具体实现可复用仓库原 pipeline 中的解析结构）。

\subsection{几何吸收（\texorpdfstring{$\rho$}{rho}）下首吸收时间生成函数}
令 target 结点为 $\mathrm{tar}$。到达 $\mathrm{tar}$ 时以概率 $\rho$ 吸收，否则以概率 $1-\rho$ 继续。设未吸收游走的访问生成函数为
\[
S_{i,j}(z)=\sum_{t\ge 0}\mathbb{P}(X_t=i\mid X_0=j)z^t.
\]
则首吸收时间 $T$ 的生成函数（$t\ge 1$ 的 pmf 系数）可写为
\begin{equation}
\tilde f(z)=\sum_{t\ge 1}f(t)z^t
=\frac{\rho\,S_{\mathrm{tar},n_0}(z)}{(1-\rho)+\rho\,S_{\mathrm{tar},\mathrm{tar}}(z)}.
\end{equation}
当 $\rho=1$ 时简化为
\[
\tilde f(z)=\frac{S_{\mathrm{tar},n_0}(z)}{S_{\mathrm{tar},\mathrm{tar}}(z)}.
\]

\subsection{Abate--Whitt (AW) 反演：用 FFT 计算 \texorpdfstring{$f(t)$}{f(t)}}
系数反演可由 Cauchy 积分与离散化得到。取 $z_m=re^{2\pi i m/L}$（$r<1$），则
\begin{equation}
f(t)\approx r^{-t}\frac{1}{L}\sum_{m=0}^{L-1}\tilde f(z_m)e^{-2\pi i mt/L},
\end{equation}
右侧为离散傅里叶反变换（iDFT），因此可用 FFT 在 $O(L\log L)$ 时间内同时得到 $t=1,\dots,T_{\max}$ 的近似 $f(t)$。仓库 pipeline 采用的 AW 参数选择（$r,L$ 的选取）与数值稳定性控制沿用原实现。

\subsection{尾部指数衰减与谱半径指标 \texorpdfstring{$\gamma$}{gamma}}
对吸收链，记 transient 子矩阵为 $Q$，则生存函数
\[
S(t)=\mathbb{P}(T>t)=\mathbf{1}^\top Q^t e_{n_0}.
\]
由于 $Q$ 为非负矩阵，且在通常连通条件下满足 Perron--Frobenius，存在最大特征值 $\lambda_{\max}\in(0,1)$ 使得
\[
S(t)\sim c\,\lambda_{\max}^t\quad (t\to\infty).
\]
因此定义尾部渐近指数衰减率
\begin{equation}
\gamma=-\log(\lambda_{\max}),
\end{equation}
作为跨 $\beta$、跨 $N$、跨 $K$ 的可比较 tail 指标。数值上我们通过构造 $Q$ 并求其谱半径得到 $\gamma$；当 $\rho=1$ 时删除 target 状态，当 $\rho<1$ 时 target 行乘 $(1-\rho)$ 表示吸收漏失。

\section{Valley 的可比较定义与窗口划分}

令首吸收时间的概率质量函数为 $f(t)=\mathbb{P}(T=t)$（$t\in\mathbb{N}^+$）。
我们采用严格局部极大作为峰的候选：
\begin{equation}
f(t)>f(t-1),\qquad f(t)>f(t+1),\qquad f(t)\ge h_{\min}.
\end{equation}
再用相对高度阈值过滤次峰：只保留 $f(t)\ge \texttt{second\_rel\_height}\cdot \max_{u} f(u)$ 的峰。

\subsection{V0：默认 Valley（两峰区间最小值）}
若检测到时间顺序上的前两个峰 $t_1<t_2$，定义
\begin{equation}
t_v=\arg\min_{t_1<t<t_2} f(t).
\end{equation}
该定义参数极少，且天然保证 $t_v$ 位于两峰之间，便于跨 $\beta$ 与跨 $K$ 的比较。

\subsection{V1：基于峰间相对距离的窗口尺度（用于轨迹分类对比）}
记峰间距 $\Delta T=t_2-t_1$，定义窗口半宽
\begin{equation}
\Delta=\max\left(1,\left\lfloor \delta_{\mathrm{frac}}\Delta T\right\rfloor\right).
\end{equation}
将每条轨迹按其吸收时间 $T$ 分配到最近的中心 $\{t_1,t_v,t_2\}$：
若 $\min_{c\in\{t_1,t_v,t_2\}} |T-c|\le \Delta$，则归入对应窗口，否则归为 other。
由于 $\Delta$ 随 $\Delta T$ 成比例缩放，该划分在不同 $\beta$ 下保持可比性。

\subsection{V2：鲁棒 Valley（排除靠近峰的区域）}
为了避免 $t_v$ 贴近某个峰（在两峰接近时可能发生），可定义
\begin{equation}
t_v=\arg\min_{t_1+\alpha\Delta T\le t\le t_2-\alpha\Delta T} f(t),
\end{equation}
其中 $\alpha\in(0,1/2)$，例如 $\alpha=0.1$。V2 作为敏感性分析口径。

\section{只改变 \texorpdfstring{$\beta$}{beta} 的控制变量扫描（\texorpdfstring{$N=100$}{N=100}）}
固定 $N=100,q=2/3,\rho=1$，口径为 \texttt{lazy\_selfloop}，仅改变 $\beta$。
图~\ref{fig:beta-sweep-key-cn} 汇总 AW 精确峰高/峰位与 tail 指标，并给出全 $\beta$ 的精确 $f(t)$ 叠加图。

\paragraph{峰-谷量化与 \texorpdfstring{$K=2$}{K=2} vs \texorpdfstring{$K=4$}{K=4}}
我们记录 valley 深度指标 $\texttt{hv\_over\_max}=h_v/\max(h_1,h_2)$ 与 $\texttt{h2\_over\_h1}$（见 CSV）。
以 $\beta=0.01$ 为例，$K=2$ 的 $\texttt{hv\_over\_max}\approx0.34$，而 $K=4$ 仅约 $0.15$，说明 $K=4$ 的 valley 更深且两峰分离更明显。
同一 $\beta$ 下 $K=4$ 的 $t_1,t_2$ 更早、$\gamma$ 更大，体现时间尺度更快、尾部更轻。
同时 $K=2$ 不存在 jump-over（$J\equiv0$），而 $K=4$ 允许 $J>0$，这会在后续 MC 轨迹分类中体现为额外的 $J$-相关类别。

\paragraph{代表性 \texorpdfstring{$f(t)$}{f(t)}（峰谷可读性增强）}
图~\ref{fig:beta-example-ft-cn} 给出 $\beta=0.01$ 的代表性 $f(t)$，时间轴聚焦在峰谷区间（避免 $t\approx0$ 的长纵轴拉伸），以突出 peak/valley 对比。

\paragraph{Tail decay 机制与诊断}
尾部渐近由 $\gamma=-\log(\lambda_{\max})$ 控制，图~\ref{fig:tail-logS-cn} 展示 $\log S(t)$ 的近线性衰减。
以 $\beta=0.01$ 为例，$K=2$ 的 $\gamma\approx3.89\times10^{-4}$，$K=4$ 为 $8.52\times10^{-4}$，与 $\gamma$-vs-$\beta$ 曲线一致。

\subsection{自动选择 \texorpdfstring{$\beta$}{beta}}
我们用统一阈值筛选（$t_2/t_1\ge10$、$\texttt{hv\_over\_max}\le0.6$、$\texttt{h2\_over\_h1}\ge0.05$），并取同时满足 $K=2$ 与 $K=4$ 的最小 $\beta$。
本次扫描得到 $\beta_\ast=0.01$（见 \texttt{outputs/selected\_beta.txt}）。

以上为精确 AW 的结果与机制解释。下一节将用 MC 轨迹分类检验 jump-over 在不同窗口的分布。

\section{MC 扫描：四分类随 \texorpdfstring{$\beta$}{beta} 变化（\texorpdfstring{$N=100$}{N=100}）}
使用 MC 统计窗口内的四类轨迹比例，重点关注 $J$-相关类别在 valley/第二峰窗口的集中程度。
窗口由 $(t_1,t_v,t_2)$ 与 $\Delta=\lfloor\delta_{\mathrm{frac}}(t_2-t_1)\rfloor$ 定义。
四类为：
$\texttt{class\_C0J0}$（不走 shortcut、无 jump-over）、
$\texttt{class\_C1pJ0}$（有 shortcut、无 jump-over）、
$\texttt{class\_C0J1p}$（无 shortcut、有 jump-over）、
$\texttt{class\_C1pJ1p}$（有 shortcut、有 jump-over）。
其中 $K=2$ 情况下恒有 $J\equiv0$，因此 $J$ 相关类别应为 0 或接近 0。
图~\ref{fig:mc-beta-classes-cn} 展示四类轨迹在不同窗口内的条件概率随 $\beta$ 的变化。

\paragraph{解释要点}
K=2 的 $J$ 类别几乎为 0，表明 jump-over 机制仅在 $K=4$ 下起作用；
在 $K=4$ 的 valley 与第二峰窗口中，$J$-相关类别（$\texttt{C0J1p}$ 与 $\texttt{C1pJ1p}$）占比明显上升，
对应了 ``跨过吸收点但未吸收'' 的时间延迟路径。

\begin{figure}[!htbp]
  \centering
  \begin{minipage}{0.32\linewidth}
    \centering
    \textbf{(a) peak1}\par
    \includegraphics[width=\linewidth]{outputs/mc_beta_sweep_N100/aux_peak1_vs_beta.png}
  \end{minipage}\hfill
  \begin{minipage}{0.32\linewidth}
    \centering
    \textbf{(b) valley}\par
    \includegraphics[width=\linewidth]{outputs/mc_beta_sweep_N100/aux_valley_vs_beta.png}
  \end{minipage}\hfill
  \begin{minipage}{0.32\linewidth}
    \centering
    \textbf{(c) peak2}\par
    \includegraphics[width=\linewidth]{outputs/mc_beta_sweep_N100/aux_peak2_vs_beta.png}
  \end{minipage}
  \caption{$\mathbb{P}(C\ge1\mid\text{window})$ 与 $\mathbb{P}(J\ge1\mid\text{window})$ 的随 $\beta$ 变化。}
  \label{fig:mc-beta-aux-cn}
\end{figure}

有了 $\beta$ 方向的轨迹结构对比，下一节固定 $\beta_\ast$ 扫 $N$，观察慢尺度随系统规模的变化。

\section{固定 \texorpdfstring{$\beta=0.01$}{beta=0.01} 改变 \texorpdfstring{$N$}{N}（精确 AW）}
取自动选择的 $\beta_\ast=0.01$，固定 $q=2/3,\rho=1$ 与 \texttt{lazy\_selfloop}，对偶数 $N$ 扫描峰高、峰位与尾部衰减指标。
图~\ref{fig:n-sweep-peaks-cn} 显示：$t_2$ 随 $N$ 增大显著右移（更慢的第二时间尺度），而 $h_2$ 逐步降低；
$\gamma$ 随 $N$ 下降，表明尾部更厚（衰减更慢），对应的曲线见图~\ref{fig:n-sweep-gamma-cn}。

以上给出精确 AW 的规模效应。下一节用 MC 轨迹分类检验 jump-over 在不同 $N$ 下的窗口占比。

\section{MC 扫描：四分类随 \texorpdfstring{$N$}{N} 变化（固定 \texorpdfstring{$\beta=0.01$}{beta=0.01}）}
保持 $\beta=0.01$，对偶数 $N$ 做 MC 扫描并统计四类轨迹在不同窗口的比例。
回顾四类定义：$C$ 表示是否至少使用一次 shortcut（$C0$ 为未用、$C1p$ 为至少一次），
$J$ 表示是否出现 jump-over（$J0$ 为无、$J1p$ 为至少一次），
因此 $\texttt{C0J0}$、$\texttt{C1pJ0}$、$\texttt{C0J1p}$、$\texttt{C1pJ1p}$ 对应四类轨迹。
图~\ref{fig:mc-n-classes-cn} 汇总了四类轨迹比例随 $N$ 的变化。

\paragraph{对比要点}
$K=2$ 中 $J$-相关类别几乎为 0；
$K=4$ 在 valley 与第二峰窗口中更容易出现 $J$-相关类别，体现 jump-over 对延迟吸收路径的贡献。

\begin{figure}[!htbp]
  \centering
  \begin{minipage}{0.32\linewidth}
    \centering
    \textbf{(a) peak1}\par
    \includegraphics[width=\linewidth]{outputs/mc_N_sweep_beta_selected/aux_peak1_vs_N.png}
  \end{minipage}\hfill
  \begin{minipage}{0.32\linewidth}
    \centering
    \textbf{(b) valley}\par
    \includegraphics[width=\linewidth]{outputs/mc_N_sweep_beta_selected/aux_valley_vs_N.png}
  \end{minipage}\hfill
  \begin{minipage}{0.32\linewidth}
    \centering
    \textbf{(c) peak2}\par
    \includegraphics[width=\linewidth]{outputs/mc_N_sweep_beta_selected/aux_peak2_vs_N.png}
  \end{minipage}
  \caption{$\mathbb{P}(C\ge1\mid\text{window})$ 与 $\mathbb{P}(J\ge1\mid\text{window})$ 随 $N$ 的变化。}
  \label{fig:mc-n-aux-cn}
\end{figure}

复现与参数细节见复现指南一节。

\section{复现指南（命令与输出）}
以下命令均在 \texttt{reports/lazy\_jump\_ext/} 内执行。

\paragraph{1. \texorpdfstring{$\beta$}{beta} 扫描（精确 AW）}
\begin{verbatim}
python code/beta_sweep_peaks_tail.py \
  --N 100 --Ks 2 4 --q 0.6666666667 --rho 1.0 --mode lazy_selfloop \
  --beta-min 0.0 --beta-max 0.2 --beta-num 21 \
  --outdir outputs/beta_sweep_N100
\end{verbatim}
输出：\texttt{outputs/beta\_sweep\_N100/beta\_sweep\_metrics\_N100.csv} 与三张曲线图。

\paragraph{2. 自动选择 \texorpdfstring{$\beta$}{beta}}
\begin{verbatim}
python code/choose_beta.py \
  --csv outputs/beta_sweep_N100/beta_sweep_metrics_N100.csv \
  --Ks 2 4 --macro_ratio 10 --hv_over_max 0.6 --min_h2_over_h1 0.05 \
  --fallback_beta 0.02 \
  --out outputs/selected_beta.txt
\end{verbatim}

\paragraph{3. 固定 \texorpdfstring{$\beta$}{beta} 扫 \texorpdfstring{$N$}{N}（精确 AW）}
\begin{verbatim}
python code/N_sweep_peaks_tail.py \
  --K 2 4 --beta 0.01 --q 0.6666666667 --rho 1.0 --mode lazy_selfloop \
  --N-min 50 --N-max 300 --N-step 2 --only-even \
  --outdir outputs/N_sweep_beta_selected
\end{verbatim}

\paragraph{4. MC 扫 \texorpdfstring{$\beta$}{beta}（四分类）}
\begin{verbatim}
python code/mc_sweep_beta_classes.py \
  --N 100 --Ks 2 4 --q 0.6666666667 --rho 1.0 --mode lazy_selfloop \
  --betas 0.001 0.005 0.01 0.02 0.03 \
  --n-walkers 200000 --seed 0 \
  --outdir outputs/mc_beta_sweep_N100
\end{verbatim}

\paragraph{5. MC 扫 \texorpdfstring{$N$}{N}（固定 \texorpdfstring{$\beta$}{beta}）}
\begin{verbatim}
python code/mc_sweep_N_classes.py \
  --beta 0.01 --Ks 2 4 --q 0.6666666667 --rho 1.0 --mode lazy_selfloop \
  --N-min 50 --N-max 300 --N-step 10 --only-even \
  --n-walkers 200000 --seed 0 \
  --outdir outputs/mc_N_sweep_beta_selected
\end{verbatim}

\paragraph{6. Tail 诊断（\texorpdfstring{$\log S(t)$}{log S(t)}）}
\begin{verbatim}
python code/tail_diagnostics.py \
  --N 100 --K 2 --beta 0.01 --q 0.6666666667 --rho 1.0 --mode lazy_selfloop \
  --max-steps-aw 12000 --outdir outputs/tail_diag
python code/tail_diagnostics.py \
  --N 100 --K 4 --beta 0.01 --q 0.6666666667 --rho 1.0 --mode lazy_selfloop \
  --max-steps-aw 12000 --outdir outputs/tail_diag
\end{verbatim}

\paragraph{7. 全 \texorpdfstring{$\beta$}{beta} 精确曲线叠加图（可选）}
\begin{verbatim}
python code/plot_aw_all_betas.py \
  --N 100 --Ks 2 4 --q 0.6666666667 --rho 1.0 --mode lazy_selfloop \
  --beta-min 0.0 --beta-max 0.2 --beta-num 21 \
  --out outputs/beta_sweep_N100/aw_all_betas_N100.png
\end{verbatim}

\paragraph{8. 编译报告}
\begin{verbatim}
latexmk -xelatex -interaction=nonstopmode ring_lazy_jump_ext_cn.tex
latexmk -pdf -interaction=nonstopmode ring_lazy_jump_ext_en.tex
\end{verbatim}


\end{document}
